%%=============================================================================
%% CAPITULO 1 -- INTRODUCAO
%% Versao publica (vitrine): apresenta o tema e o contexto, sem a pergunta de
%% pesquisa, as hipoteses e o desenho metodologico detalhado.
%%=============================================================================

\section{Introdução}

A incorporação de tecnologias de automação e de inteligência artificial (IA)
tem ganhado espaço nas estratégias de modernização do Estado brasileiro,
impulsionada pelo crescimento do volume de dados e pela exigência
constitucional de eficiência administrativa \parencite[art.~37,
caput]{brasil1988}. Nesse contexto, a Advocacia Pública, função essencial à
Justiça responsável pela representação judicial e pela consultoria jurídica
dos entes federativos \parencite[arts.~131 e 132]{brasil1988}, enfrenta
desafios particulares relacionados ao desempenho de suas atribuições. No
âmbito da representação judicial, que constitui o objeto deste estudo,
destacam-se a elevada demanda processual e o expressivo volume de intimações
recebidas, cuja gestão envolve não apenas a análise jurídica, mas também um
conjunto de atividades padronizadas de identificação, classificação, registro
e encaminhamento.

O problema é visível na triagem de intimações. Em demandas repetitivas e de
baixa complexidade decisória, a definição da providência cabível tende a ser
rápida para o profissional experiente; o custo operacional concentra-se menos
nesse juízo do que no tratamento individual de cada ocorrência: abrir,
consultar, conferir, registrar e encaminhar, sucessivamente, ao longo da
jornada. A Figura~\ref{fig:ciclo-intimacao} representa esse ciclo. Nesse
fluxo, a IA pode auxiliar na classificação da intimação e na sugestão das
providências cabíveis, enquanto mecanismos de automação podem assumir as
etapas operacionais delas decorrentes.

\begin{figure}[htb]
  \caption{Etapas do tratamento individual de uma intimação}
  \label{fig:ciclo-intimacao}
  \centering
  \begin{tikzpicture}[
      etapa/.style={draw, rounded corners, align=center, minimum height=1.1cm,
                    text width=2.1cm, font=\small},
      seta/.style={-{Stealth[length=2.5mm]}, thick},
      node distance=0.45cm]
    \node[etapa] (abr) {Abrir};
    \node[etapa, right=of abr] (con) {Consultar};
    \node[etapa, right=of con] (cnf) {Conferir};
    \node[etapa, right=of cnf] (reg) {Registrar};
    \node[etapa, right=of reg] (enc) {Encaminhar};
    \draw[seta] (abr) -- (con);
    \draw[seta] (con) -- (cnf);
    \draw[seta] (cnf) -- (reg);
    \draw[seta] (reg) -- (enc);
  \end{tikzpicture}
  \fonte{Elaborada pela autora (2026).}
\end{figure}

O Núcleo de Cumprimento de Sentença (NCS) da Procuradoria-Geral do Distrito
Federal (PGDF), unidade selecionada como estudo de caso, é especializado em
cumprimentos individuais decorrentes de julgados coletivos e parece reunir
características adequadas a essa investigação.
